\uuid{m2vU}
\titre{Variation de la constante/prolongement par continuité}
\auteur{hanani ; mzari}
\organisation{exo7}
\licence{CC BY-NC-SA 4.0}
\datecreate{2020-01-20}
\niveau{L1}
\module{Analyse}
\chapitre{Equation différentielle}
\sousChapitre{Résolution d'équation différentielle du premier ordre}
\difficulte{3}
\theme{Equations du premier ordre, variation de la constante, prolongement par continuité}
\exo7qcmid{lille-2.equations-differentielles.015}

\contenu{
\begin{question}
\qclassification{225.01}
\qtags{motcle=variation de la constante/prolongement par continuité}
On considère l'équation différentielle $(E)$ : $\displaystyle y'\ln x+\frac{y}{x}=2x$ sur $]1,+\infty[$. Parmi les affirmations suivantes cocher celles qui sont vraies :
\begin{answers}
\bad{La solution générale de l'équation homogène est : $\displaystyle y(x)=k\ln x$, $k\in \Rr$.}
\good{Si $\displaystyle y(x)=\frac{k(x)}{\ln x}$ est une solution de $(E)$ sur $]1,+\infty[$, alors $k'(x)=2x$.}
\good{La solution générale de $(E)$ sur $]1,+\infty[$ est $\displaystyle y(x)=\frac{k+x^2}{\ln x}$, $k\in \Rr$.}
\bad{$(E)$ possède une infinité de solutions sur $[1,+\infty[$.}
\end{answers}
\begin{explanations}
Une primitive de $\displaystyle \frac{1}{x\ln x}$ est $\ln |\ln x|$, donc la solution générale de l'équation homogène est $\displaystyle y(x)=\frac{k}{\ln x}$. Si $\displaystyle y(x)=\frac{k(x)}{\ln x}$ est une solution de $(E)$ sur $]1,+\infty[$, alors $k'(x)=2x$ et donc $k=x^2+C$. Ainsi la solution générale de $(E)$ sur $]1,+\infty[$ est : $\displaystyle y(x)=\frac{k+x^2}{\ln x}$, $k\in \Rr$. On pose $x=1$ dans $(E)$, on aura $y(1)=2$. Or $\displaystyle \lim _{x\to 1^+}\frac{k+x^2}{\ln x}=2$ si, et seulement si, $k=-1$. Donc $(E)$ ne peut admettre une infinité de solutions sur $[1,+\infty[$.
\end{explanations}
\end{question}
}
